
\subsubsection{A Welch/Frobenius Upper Bound in the Isotropic Key/Value Regime}
\phantomsection\label{app:theory:margin_bounds:rkrv:welch}
\label{app:welch-upper-bound-iso}

In the main text, \Cref{thm:bilinear-scaling} provides a lower bound on the minimum margin achieved by our construction in the isotropic key/value regime.
We now show that this lower bound is asymptotically tight, up to constants and logarithmic factors, by proving a matching \emph{upper bound} on the best possible margin of any admissible rank-limited kernel memory in this regime. The key idea is that the arbitrary-keys / isotropic-values analysis in \Cref{app:theory:margin_bounds:akrv} is controlled by the column energy
\[
E_{\mathrm{col}}(i) := \sum_{t\neq i} \hat{\matK}_{ti}^2.
\]
Thus, to upper bound the achievable margin, it suffices to show that for any rank-$m$ kernel, some column must have nontrivial off-diagonal squared mass.
We achieve this using a Welch / Frobenius argument: intuitively, low rank forces a non-negligible amount of kernel mass off the diagonal, which in turn yields an unavoidable cross-talk floor.

We begin by formalizing the class of kernels under consideration.

\begin{definition}[Rank-$m$ kernel]
\label{def:rank-m-kernel}
Let $\vk_{1},\dots,\vk_{F} \in \mathbb{R}^d$ be the stored keys. We say that
$\hat{\matK} \in \mathbb{R}^{F\times F}$ is a \emph{rank-$m$ kernel} if there exists a feature
map $\phi(\cdot)$ with $\phi(\vk_{i})\in\mathbb{R}^m$ such that
\[
\hat{\matK}_{ij} = \langle \phi(\vk_{i}), \phi(\vk_{j})\rangle
\qquad\text{for all } i,j\in[F].
\]
Equivalently, $\hat{\matK} \succeq 0$ and $\operatorname{rank}(\hat{\matK})\le m$.
\end{definition}

The following lemma is the kernel-side ingredient. It lower bounds the maximum column energy
of any PSD rank-$m$ kernel in terms of its diagonal.

\begin{lemma}[Welch/Frobenius lower bound for rank-$m$ kernels]
\label{lem:welch-frobenius-general}
Let $\hat{\matK} \in \mathbb{R}^{F\times F}$ be a rank-$m$ kernel in the sense of
Definition~\ref{def:rank-m-kernel}. Define
\[
E_{\mathrm{col}}(i) := \sum_{t\neq i}\hat{\matK}_{ti}^2,
\qquad
E_{\mathrm{col}} := \max_{i\in[F]} E_{\mathrm{col}}(i).
\]
Then
\[
\|\hat{\matK}\|_{F}^2 \ge \frac{\operatorname{tr}(\hat{\matK})^2}{m},
\]
and consequently
\[
E_{\mathrm{col}}
\;\ge\;
\frac{1}{F}\left(
\frac{\operatorname{tr}(\hat{\matK})^2}{m}
-
\sum_{i=1}^F \hat{\matK}_{ii}^2
\right).
\]
In particular, if
\[
\hat{\matK}_{ii}\in[1-\varepsilon_{\mathrm{diag}},\,1+\varepsilon_{\mathrm{diag}}]
\qquad\text{for all } i\in[F],
\]
then
\[
E_{\mathrm{col}}
\;\ge\;
\frac{F(1-\varepsilon_{\mathrm{diag}})^2}{m}
-
(1+\varepsilon_{\mathrm{diag}})^2.
\]
\end{lemma}

\begin{proof}
Since $\hat{\matK} \succeq 0$ and $\operatorname{rank}(\hat{\matK})\le m$, let
$\lambda_{1},\dots,\lambda_{r}$ be the nonzero eigenvalues of $\hat{\matK}$, where $r\le m$. Then
\[
\operatorname{tr}(\hat{\matK})=\sum_{a=1}^r \lambda_{a},
\qquad
\|\hat{\matK}\|_{F}^2=\sum_{a=1}^r \lambda_{a}^2.
\]
Applying Cauchy--Schwarz to $(\lambda_{1},\dots,\lambda_{r})$ gives
\[
\Big(\sum_{a=1}^r \lambda_{a}\Big)^2
\le r\sum_{a=1}^r \lambda_{a}^2
\le m\sum_{a=1}^r \lambda_{a}^2,
\]
hence
\[
\|\hat{\matK}\|_{F}^2 \ge \frac{\operatorname{tr}(\hat{\matK})^2}{m}.
\]

Next expand the Frobenius norm entrywise:
\[
\|\hat{\matK}\|_{F}^2
=
\sum_{i=1}^F \hat{\matK}_{ii}^2
+
\sum_{i=1}^F\sum_{t\neq i}\hat{\matK}_{ti}^2
=
\sum_{i=1}^F \hat{\matK}_{ii}^2
+
\sum_{i=1}^F E_{\mathrm{col}}(i).
\]
Therefore
\[
\sum_{i=1}^F E_{\mathrm{col}}(i)
\ge
\frac{\operatorname{tr}(\hat{\matK})^2}{m}
-
\sum_{i=1}^F \hat{\matK}_{ii}^2.
\]
Since the maximum is at least the average,
\[
E_{\mathrm{col}}
=
\max_{i} E_{\mathrm{col}}(i)
\ge
\frac{1}{F}\sum_{i=1}^F E_{\mathrm{col}}(i)
\ge
\frac{1}{F}\left(
\frac{\operatorname{tr}(\hat{\matK})^2}{m}
-
\sum_{i=1}^F \hat{\matK}_{ii}^2
\right).
\]

If moreover $\hat{\matK}_{ii}\in[1-\varepsilon_{\mathrm{diag}},1+\varepsilon_{\mathrm{diag}}]$ for all $i$,
then
\[
\operatorname{tr}(\hat{\matK})=\sum_{i=1}^F \hat{\matK}_{ii} \ge F(1-\varepsilon_{\mathrm{diag}})
\]
and
\[
\sum_{i=1}^F \hat{\matK}_{ii}^2 \le F(1+\varepsilon_{\mathrm{diag}})^2.
\]
Substituting these into the previous display yields
\[
E_{\mathrm{col}}
\ge
\frac{1}{F}\left(
\frac{F^2(1-\varepsilon_{\mathrm{diag}})^2}{m}
-
F(1+\varepsilon_{\mathrm{diag}})^2
\right)
=
\frac{F(1-\varepsilon_{\mathrm{diag}})^2}{m}
-
(1+\varepsilon_{\mathrm{diag}})^2.
\]
\end{proof}

Regarding the assumption
\[
\hat{\matK}_{ii}\in[1-\varepsilon_{\mathrm{diag}},1+\varepsilon_{\mathrm{diag}}],
\]
note that for the exact quadratic kernel on unit-norm isotropic keys, one has
\[
K_{2}(\vk_{i},\vk_{i})=\langle \vk_{i},\vk_{i}\rangle^2 = 1
\]
exactly.
For the sketched / bilinear random-feature kernel used in our construction, the diagonal is generally not exactly $1$, but it instead concentrates near $1$ with high probability by \Cref{lem:Kdiag-lower-bilinear}.

Combining the kernel-side lower bound with the isotropic-value competitor argument used in the
proof of \Cref{rem:welch-benchmark} yields the desired asymptotic optimality statement.

\begin{corollary}[Rank-$m$ kernels are asymptotically unbeatable in the isotropic key/value regime]
\label{cor:rank-m-unbeatable-iso}
Assume the isotropic key/value setting, and let $\hat{\matK}$ be any rank-$m$ kernel independent of
the values, with
\[
\hat{\matK}_{ii}\in[1-\varepsilon_{\mathrm{diag}},\,1+\varepsilon_{\mathrm{diag}}]
\qquad\text{for all } i\in[F].
\]
Then the isotropic-value competitor argument used in the proof of \Cref{rem:welch-benchmark} yields
\[
\gamma_{\min}
\;\le\;
(1+\varepsilon_{\mathrm{diag}})
-
\Omega\!\left(
\sqrt{\frac{E_{\mathrm{col}}\log F}{d}}
\right)
+
\mathcal{O}\!\left(
\sqrt{\frac{\log F}{d}}
\right)
+
\mathcal{O}\!\left(
\sqrt{\frac{E_{\mathrm{col}}\log(1/\delta)}{d}}
\right)
+
\mathcal{O}\!\left(
\sqrt{\frac{E_{\mathrm{col}}}{d}}
\right).
\]
Combining with Lemma~\ref{lem:welch-frobenius-general} yields
\[
\gamma_{\min}
\;\le\;
(1+\varepsilon_{\mathrm{diag}})
-
\Omega\!\left(
\sqrt{
\frac{\big(\frac{F(1-\varepsilon_{\mathrm{diag}})^2}{m}-(1+\varepsilon_{\mathrm{diag}})^2\big)\log F}{d}
}
\right)
+ \text{lower-order terms}.
\]
In particular, if $\varepsilon_{\mathrm{diag}}=o(1)$ and $F/m\to\infty$, then up to constants and
logarithmic factors,
\[
\gamma_{\min}
\;\le\;
1 - \Omega\!\left(\sqrt{\frac{(F/m)\log F}{d}}\right).
\]
Equivalently, no admissible rank-$m$ kernel can asymptotically beat the
$\sqrt{(F/m)\log F / d}$ margin floor in the isotropic key/value regime, up to constants and logs.
\end{corollary}

\begin{proof}
Apply the isotropic-value competitor argument used in the proof of \Cref{rem:welch-benchmark} to
the column $i^\star$ with $E_{\mathrm{col}}(i^\star)=E_{\mathrm{col}}$. This gives
\[
\gamma_{\min}
\le
\hat{\matK}_{i^\star i^\star}
-
\Omega\!\left(
\sqrt{\frac{E_{\mathrm{col}}\log F}{d}}
\right)
+
\mathcal{O}\!\left(
|\hat{\matK}_{i^\star i^\star}|\sqrt{\frac{\log F}{d}}
\right)
+
\mathcal{O}\!\left(
\sqrt{\frac{E_{\mathrm{col}}\log(1/\delta)}{d}}
\right)
+
\mathcal{O}\!\left(
\sqrt{\frac{E_{\mathrm{col}}}{d}}
\right),
\]
after absorbing harmless constant-factor differences into the big-$\mathcal{O}$/big-$\Omega$
notation.

Using $\hat{\matK}_{i^\star i^\star}\le 1+\varepsilon_{\mathrm{diag}}$ gives the first bound.
The second bound results from substituting the lower bound on $E_{\mathrm{col}}$ from~\Cref{lem:welch-frobenius-general}.

Finally, if $\varepsilon_{\mathrm{diag}}=o(1)$, then
\[
\frac{F(1-\varepsilon_{\mathrm{diag}})^2}{m}
-
(1+\varepsilon_{\mathrm{diag}})^2
=
\Theta\!\left(\frac{F}{m}\right)
\]
whenever $F/m\to\infty$, so the dominant negative term scales as
\[
\sqrt{\frac{(F/m)\log F}{d}},
\]
which proves the last claim.
\end{proof}
